\section{Considerações Parciais}

Este trabalho apresentou a proposta de uma plataforma web que junta um fórum
acadêmico organizado por disciplina e um sistema de voluntariado voltado para a
comunidade da UNIFEI. Foi feito o levantamento bibliográfico nos eixos de
comunicação em ambientes educacionais, fóruns de discussão, aprendizagem
colaborativa e voluntariado universitário, junto com a análise comparativa entre
as plataformas usadas pela UNIFEI, USP, UFMG e UNESPAR, que acabou deixando
claro algo que se repete em todas as instituições analisadas: a falta de uma
ferramenta que junte a discussão acadêmica estruturada com mecanismos de
engajamento e com uma gestão integrada do voluntariado.

A partir desse levantamento, foram documentados 53 requisitos funcionais e mais
de 30 requisitos não funcionais, foi feita a modelagem do banco de dados com 17
tabelas em terceira forma normal, foi definida a arquitetura do sistema com a
separação entre frontend, backend, banco de dados e cache, e foram construídos
os diagramas de entidade-relacionamento, casos de uso, processos e implantação.
O escopo do MLP foi delimitado dando prioridade às funcionalidades que entregam
valor e uma boa experiência ao usuário desde a primeira versão, contemplando o
login por CPF, a criação de tópicos ligados às disciplinas, a votação nas
respostas, a marcação de melhor resposta, a busca textual e, no módulo de
voluntariado, a inscrição nas oportunidades com a certificação integrada.

Para o TCC-2 estão previstas a implementação do protótipo funcional, usando o
React.js com TypeScript no frontend e o Django com o Django REST Framework e o
Django Channels no backend, e a realização de testes com usuários da UNIFEI
para validar se a plataforma de fato responde aos problemas que foram
levantados. Também está prevista a integração de um assistente inteligente ao
fórum, que vai reconhecer o que o aluno está procurando e direcionar a sua
dúvida para o lugar certo, seja o contato de um professor, uma informação da
UNIFEI, o material de uma matéria, um tópico parecido que já foi respondido
antes, a disciplina correta ou uma oportunidade de voluntariado, ajudando o
aluno a chegar mais rápido no que precisa sem se perder no caminho. Além disso,
está previsto o desenvolvimento das funcionalidades que ficaram planejadas para
as próximas versões, como as notificações em tempo real, o ranking semestral, a
moderação avançada e a integração com o SIGAA, deixando o sistema cada vez mais
perto de atender o que a comunidade acadêmica realmente precisa no dia a dia.